\chapter{Reasonable spaces}\label{chap:Reasonable spaces}

\begin{wrapfigure}{r}{23 mm}
\vskip-8mm
\centering
\begin{tikzcd}
& \c{Y}\arrow[d, "w"]\\
\tilde{\c{Z}} \arrow[r, "u"]\arrow[ur, "v"] & \c{Z}
\end{tikzcd}
\end{wrapfigure}

Most of connected spaces $\c{Z}$ that appear in applications admit a covering $u\:\tilde{\c{Z}}\to \c{Z}$ with a simply-connected total space $\tilde{\c{Z}}$ and such covering has a universal property --- any other connected covering $w\:\c{Y}\to \c{Z}$ can be included in the commutative diagram of coverings.
By that reason, $u$ is called a \textit{universal covering} of~$\c{Z}$.

We already proved several statement, including \ref{thm:bijection-fiber} and \ref{ex:monodromy:isomorphism}, showing that a universal covering knows a lot about fundamental group of the base.
In this chapter we discuss spaces for which the last statement fails to be true and introduce a mild condition that will be latter used to prove existence of a universal covering.

\section{Definitions}\label{sec:Reasonable spaces}



Consider a connected covering $u\:\tilde{\c{Z}}\to \c{Z}$.
This covering will be called \emph{universal} if any other connected covering $w\:\c{Y}\to \c{Z}$ can be included in the commutative diagram of coverings.

In the next chapter we will show that every reasonable connected space $\c{Z}$ admit a universal covering;
moreover, a covering $u\:\tilde{\c{Z}}\to \c{Z}$ is universal if and only if the total space $\tilde{\c{Z}}$ is simply-connected.

The term reasonable will be used as a shortcut for \index{locally simply-connected}\emph{locally simply-connected}.
This term will be used only in these notes; it is not standard and can mean something else elswere.

So, a topological space $\c{Z}$ is \index{reasonable space}\emph{reasonable} if for any open set $V\subset \c{Z}$ and any $p\in V$ there is an open simply-connected set $W$ such that $V\z\supset W\ni p$.

For instance, Euclidean space is reasonable;
it follows since any open ball in a Euclidean space is convex; in particular, contractible and simply-connected.

Since any simply-connected set is path-connected, any reasonable space is locally path-connected and locally connected;
see \ref{sec:Locally path-connected spaces} and \ref{sec:Locally connected spaces}.
In particular, from \ref{obs:open-connected-component} we have that every connected component of reasonable space is clopen.


\begin{thm}{Exercise}\label{ex:reasonable}
Let $w\:\c{Y}\to \c{Z}$ be a covering.
Show the following:
\begin{subthm}{}
If $\c{Z}$ is reasonable then so is $\c{Y}$, and the restriction of $w$ to any connected component of $\c{Y}$ is a covering.
\end{subthm}

\begin{subthm}{}
If $w$ is surjective and $\c{Y}$ is reasonable, then so is $\c{Z}$.
\end{subthm}

\end{thm}

\section{Hawaiian earring and its relatives}

\begin{wrapfigure}[6]{r}{27 mm}
\vskip-3mm
\centering
\includegraphics{mppics/pic-1705}
\end{wrapfigure}

Let us denote by $\c{H}$ the \index{Hawaiian earring}\emph{Hawaiian earring}; it is a topological space homeomorphic to the union of nested sequence of circles in the plane shown in the picture;
the radii of circles converge to zero and the circles have one common point.


\begin{thm}{Claim}\label{clm:earring}
The Hawaiian earring $\c{H}$ is not reasonable;
it is a path-connected, locally path-conneccted, but not locally simply-connected.
\end{thm}

\parit{Proof.}
It is straightforward to check path-connectedness and local path-connecctedess of $\c{H}$.
Let us show that $\c{H}$ is not locally simply-connected.


Let $p$ be the common point of the circles in $\c{H}$.
An arbitrarily negibohood $V$ of $p$ contais a circle, say $S$.
Consider the map $r\:V\to S$ that is identity on $S$ and maps $V\setminus S$ to $p$.
Note that $r$ is a retraction.

Recall that $S$ has nontrival fundamental group; see \ref{thm:pi1(S1)}.
By \ref{ex:retract*:0},  $V$ is not simply-connected.
Hence the statement follows.
\qeds

\begin{thm}{Exercise}\label{ex:Hawaiian/bouquet}
Let $\c{H}$ be the Hawaiian earring, and let $\c{F}_\infty$ be the bouquet of countably many circles;
that is, a topological space obtained by gluing together a sequence of circles along a single point.
In other words,
\[\c{F}_\infty=(\SS^1\times \NN)/(\{p\}\times \NN),\]
where $\NN$ denotes the set of positive integers.

Construct a continuus bijection $\c{F}_\infty\to\c{H}$.
Show that $\c{F}_\infty\not\simeq \c{H}$.

\end{thm}



{

\begin{wrapfigure}{r}{23 mm}
\vskip-4mm
\centering
\begin{tikzcd}
& \c{Y}\arrow[d, "w"]\\
\c{X} \arrow[r, "u"]\arrow[ur, dotted, "\nexists" description] & \c{H}
\end{tikzcd}
\end{wrapfigure}

\begin{thm}{Exercise}\label{ex:Hawaiian-nouni}
Show that Hawaiian earring $\c{H}$ does not admit a universal covering;
that is, for any a covering $w\:\c{X}\to\c{H}$ with connected $\c{X}$ there is another covering $v\:\c{Y}\to\c{H}$ with connected $\c{Y}$
that cannot be included in the commutative diagram.
\end{thm}

}

The following exercise implies that composition of covering might fail to be a covering; according to \ref{ex:reasonable:u}, reasonable spaces behave better.

\begin{thm}{Advanced exercise}\label{ex:Hawaiian-comp-cover}
Construct a $\ZZ$-covering of the Hawaiian earring $w\:\c{Y}\to \c{H}$ and $\ZZ_2$-covering $v\:\c{X}\to \c{Y}$ such that the composition $v\circ w\:\c{X}\to \c{H}$ is not a covering.
\end{thm}

The cone over the Hawaiian earring provides an example of a contractible space (in particular, simply-connected) that is not locally simply-connected.
To construct this space think that $\c{H}$ lie in the $(x,y)$-plane in $\RR^3$;
choose a tip of the cone, say $q$, under the $(x,y)$-plane
and consider the union of line segments from $q$ to the points in~$\c{H}$.

\begin{figure}[h!]
\begin{minipage}{.38\textwidth}
\centering
\includegraphics{mppics/pic-1708}
\end{minipage}\hfill
\begin{minipage}{.58\textwidth}
\centering
\includegraphics{mppics/pic-1710}
\end{minipage}
\end{figure}

The so-called \index{Griffiths space}\emph{Griffiths space} $\c{G}$ (right picture) is obtained by gluing together two copies of such a cone at the common point of the circles in $\c{H}$.
It gives an example of a space that is path-connected, locally path-connected, and not simply-connected, but does not admit a nontrivial covering.
So this space admits trivial universal covering $\id_{\c{G}}\:\c{G}\to \c{G}$, but its total space is not simply-connected.

\begin{thm}{Very advanced exercise}\label{ex:griffiths}
Show that the Griffiths space has nontrivial fundamental group.
\end{thm}

\section{Polish circle}

\emph{Polish circle} $\c{P}$ is an example of opological space shown on the right side of the picture.
It is is a simply-conneced space that is not locally path-connected.
On the left side you see a total space of its 2-fold cover.
\begin{figure}[h!]
\vskip-0mm
\centering
\includegraphics{mppics/pic-1715}
\end{figure}
This simply-connected space admits nontrivial coverings.
The total space of the trivial covering $\id_{\c{P}}\:\c{P}\to\c{P}$ is simply-connected, but this covering does not meet the universal property discussed in \ref{sec:Reasonable spaces}.
In fact $\c{P}$ admits a universal covering with the total space shown below, but this space is not path-connected and therefore not simply-connected.
\begin{figure}[h!]
\vskip-0mm
\centering
\hspace*{-40mm}\includegraphics{mppics/pic-1720}
\end{figure}

\section{Reasonable coverings}

From now on we restrict attentions to reasonable spaces.
The following claim is the main additional feature of coverings of reasonable spaces.
Note that it does not hold for the polish circle, so the claim does not hold without assuming that reasonability.
(This time local path-connectedness will suffice.)

\begin{thm}{Claim}\label{clm:sc-base}
Let $w\:\tilde{\c{Z}}\z\to \c{Z}$ be a covering map.
Assume $\c{Z}$ is reasonable simply-connected space then $w$ is tivial.
\end{thm}

\parit{Proof.}
By \ref{ex:reasonable}, $\tilde{\c{Z}}$ is resonable, so each connected component of $\tilde{\c{Z}}$ is clopen, and
it is sufficient to show that restriction of $w$ to any connected component of $\tilde{\c{Z}}$ is a homeomorphism.

In other words, we need to show that $w$ is a homeomorphism, assuming that $\tilde{\c{Z}}$ is connected.

By \ref{thm:covering-open}, $w$ is open.
By \ref{ex:covering-clopen}, $w(\tilde{\c{Z}})$ is a clopen subset of $\c{Z}$.
Since $\c{Z}$ is simply-connected, $w$ is surjective.

It remains to show that $w$ is injective;
indeed, an open bijective continuous map must be a homeomorphism.
This follows from \ref{ex:induced-hom-covering:lift}, but let us prove it directly.

Assume the contrary, then there are two points $\tilde z_0,\tilde z_1\in\tilde{\c{Z}}$ that map to a single point, say $z_0=w(\tilde z_0)=w(\tilde z_1)$.
Since $\tilde{\c{Z}}$ is  path-connected, we can choose a path $\tilde a$ from $\tilde z_0$ to $\tilde z_1$.
Note that $a=w\circ \tilde a$ is a loop in $\c{Z}$.
Since $\c{Z}$ is simply-connected $a\sim e_{z_0}$.
Now, \ref{lem:path-lifting} implies that $\tilde z_0=\tilde a(0)=\tilde a(1)=\tilde z_1$ --- a contradiction.
\qeds

{

\begin{wrapfigure}{r}{33 mm}
\vskip-6mm
\centering
\begin{tikzcd}
& \c{Y}\arrow[dr, "w"]&\\
\c{X} \arrow[rr, "u"]\arrow[ur, "v"] && \c{Z}
\end{tikzcd}
\end{wrapfigure}

\begin{thm}{Exercise}\label{ex:reasonable-diagram}
Assume that each map in the commutative diagram is continuous and the space $\c{Z}$ is reasonable.
Show that
\begin{subthm}{ex:reasonable:u}
if $v$ and $w$ are coverings, then so is $u$,
\end{subthm}

\begin{subthm}{ex:reasonable:v}
if $u$ and $w$ are coverings, then so is  $v$,
\end{subthm}

\begin{subthm}{ex:reasonable:w}
if $u$ and $v$ are coverings, and $v$ is surjective, then $w$ is a covering.
\end{subthm}

\end{thm}

}

\begin{thm}{Exercise}\label{ex:vankam-nonconnect}
Let $V$ and $W$ be two connected open sets in a reasonable space $\c{Z}$.
Assume $\c{Z}=V\cup W$ and the intersection $V\cap W$ has at least two connected components.
Glue a nontrivial 2-fold cover $\tilde{\c{Z}}$ of $\c{Z}$ from two copies of $V$ and two copies of $W$.
Conclude that $\c{Z}$ is not simply-connected.
\end{thm}

\section{Lifts again}

\begin{thm}{Theorem}\label{thm:lifting-criterion}
Let $w\:(\tilde{\c{Z}},\tilde z_0)\z\to (\c{Z},z_0)$ be a pointed covering, and let $f\:(\c{Y},y_0)\z\to (\c{Z},z_0)$ be continuous pointed map.
Then there exists a unique pointed lift $g\:(\c{Y},y_0)\to (\tilde{\c{Z}},\tilde z_0)$.
\end{thm}

{

\begin{wrapfigure}{r}{33 mm}
\vskip-9mm
\centering
\begin{tikzcd}
& (\tilde{\c{Z}},\tilde z_0)\arrow[d, "w"]\\
{(\c{Y},y_0)} \arrow[r, "f"']\arrow[ur, "g", dashed] & (\c{Z},z_0)
\end{tikzcd}
\end{wrapfigure}

Before reading the proof, refresh the definition of fibered product from \ref{sec:Fibered product}.


\parit{Proof.} The uniqueness follows from \ref{ex:unique-lift}.

}

{

\begin{wrapfigure}{l}{33 mm}
\vskip-3mm
\centering
\begin{tikzcd}
(\tilde{\c{Y}},\tilde y_0)\arrow[d, "\tilde w"'] \arrow[r, "\tilde f"] & (\tilde{\c{Z}},\tilde z_0) \arrow[d, "w"]\\
(\c{Y},y_0) \arrow[r, "f"'] \arrow[u, bend right=50, "u"', dotted]
 & (\c{Z},z_0)
\end{tikzcd}
\end{wrapfigure}

Consider the fibered product of $f$ and $w$, described by the diagram.
By \ref{ex:pull-back}, $\tilde w\:\tilde{\c{Y}}\z\to \c{Y}$ is a covering.
Since $\c{Y}$ is reasonable and simply-connected, by \ref{clm:sc-base}, $\tilde w$ is tivial.
In particular, there is a continuous pointed map $u\:(\c{Y},y_0)\to (\tilde{\c{Y}},\tilde y_0)$ such that the diagram mostly commutes except $u\circ \tilde w\ne \id_{\tilde{\c{Y}}}$.

}

It remains to observe that $g=\tilde f\circ u$ is a requred lift of $f$.
\qeds

\begin{wrapfigure}[4]{r}{23 mm}
\vskip-5mm
\centering
\begin{tikzcd}[column sep={1em}, row sep={2em}]
\c{X}\arrow[dr, "v"']\arrow[rr, "h", dashed,leftrightarrow]& &\c{Y}\arrow[dl, "w"]\\
&\c{Z}&
\end{tikzcd}
\end{wrapfigure}

\begin{thm}{Exercise}\label{ex:ucover-unique}
Let $v\:(\c{X},x_0)\z\to (\c{Z},z_0)$ and $w\:(\c{Y},y_0)\z\to (\c{Z},z_0)$ be a coverings.
Suppose $\c{Z}$ is a connected reasonable space and both spaces $\c{X}$ and $\c{Y}$ are simply-connected.
Show that there is a unique homeomorphism $h\:(\c{X},x_0)\leftrightarrow(\c{Y},y_0)$ that makes diagram commute.
\end{thm}

\chapter{Universal covering}\label{chap:unii-covering}

\section{Existence}

\begin{thm}{Theorem}\label{thm:u-cover-exist}
Any connected reasonable space $\c{Z}$ admits a covering $u\:\tilde{\c{Z}}\to\c{Z}$ with simply-connected total space $\tilde{\c{Z}}$.
\end{thm}

\parit{Sketch of proof.}
We will construct a \textit{set} $\tilde{\c{Z}}$ with a map $u\:\tilde{\c{Z}}\to\c{Z}$, then equip $\tilde{\c{Z}}$ with a topology and check that $u$ meets the conditions in the theorem.

\parit{The set and the map.}
Choose a point $z_0\in \c{Z}$.
Denote by $\Path(\c{Z},z_0)$ the set of all paths that stat at $z_0$;
that is, a path $a\:[0,1]\to \c{Z}$ belongs to $\Path(\c{Z},z_0)$ if $a(0)=z_0$.

Let $\tilde{\c{Z}}$ be the set of all homotopy classes in $\Path(\c{Z},z_0)$; that is,
\[\tilde{\c{Z}}=\Path(\c{Z},z_0)/\sim,\]
where $\sim$ stands for the equivalence relation defined by the homotopy of paths;
see \ref{sec:Homotopy of paths}.
Recall that if $a_0\sim a_1$ for $a_0, a_1\z\in \Path(\c{Z},z_0)$, then
$a_0(1)=a_1(1)$.
Therefore, $u\:[a]\mapsto a(1)$ is a well-defined map $u\:\tilde{\c{Z}}\to\c{Z}$.


\parit{The topology.}
So far $\tilde{\c{Z}}$ is just a set; this step equips it with topology.

Given an open simply-connected set $V\subset \c{Z}$ and a path $a\in\Path(\c{Z},z_0)$ that ends in $V$,
consider the set $\tilde V_{[a]}\subset \tilde{\c{Z}}$ defined by
\[\tilde V_{[a]}=\set{[a*b]\in\tilde{\c{Z}}}{b\ \text{is a path in}\ V\ \text{that starts at}\ a(1)}.\]
Note that $\tilde V_{[a]}$ depends only on the homotopy class of $a$.

The required topology on $\tilde{\c{Z}}$ can be described by the following property:
for any open simply-connected set $V\subset \c{Z}$, and any path $a$ from $z_0$ to a point in $V$,
the set $\tilde V_{[a]}$ is open and $u$ defines a homeomorphisms $\tilde V_{[a]}\to V$.%
\footnote{In this case the sets $\tilde V_{[a]}$ for every open simply-connected set $V\subset \c{Z}$ and every path $a\in\Path(\c{Z},z_0)$ that ends in $V$ form a base of topology of $\tilde{\c{Z}}$.}
Checking that this property uniquely defines a topology is generously left to the reader.

If $a$ and $b$ be two paths from $z_0$ to a point in $V$, then
\[\tilde V_{[a]}\cap\tilde V_{[b]}\ne\emptyset
\quad\Rightarrow\quad
\tilde V_{[a]}=\tilde V_{[b]};\]
again, the proof is left to the reader.
In particular, $V$ is evenly covered by $u$ and the sets $\tilde V_{[a]}$ are the slices over $V$.
Since $\c{Z}$ is reasonable, this implies that $u$ is a covering.

Let $\tilde z_0=[e_{z_0}]$; that is, the point $\tilde z_0\in \tilde{\c{Z}}$ corresponds to the constant path $e_{z_0}$ in $\c{Z}$.
Note that $u(\tilde z_0)=e_{z_0}(1)=z_0$.

Given a path $a$, and $s\in [0,1]$, let $a_s\:[0,1]\to \c{Z}$ be the path defined by $a_s(t)=a(s\cdot t)$;
note that $a_s$ runs from $a(0)$ to $a(s)$.%
\footnote{Note that $a_s$ is not a homotopy of paths, but it is a free homotopy; see \ref{sec:Homotopy of paths}.}

Let $\tilde z=[a]\in \tilde{\c{Z}}$ for some $a\in\Path(\c{Z},z_0)$.
By \ref{lem:path-lifting}, there is a unique lift $\tilde a$ of $a$ such that $\tilde a(0)=\tilde z_0$.
Observe that $\tilde a(s)\z= [a_s]$ for any $s\in[0,1]$.
In particular, $\tilde a$ is a path from $\tilde z_0$ to $\tilde z$;
since $\tilde z$ is arbitrary, we get that $\tilde{\c{Z}}$ is path-connected.

Choose a loop $\tilde a\:[0,1]\to  \tilde{\c{Z}}$ based $\tilde z_0$.
Let $a=u\circ \tilde a$; it is a loop based at $z_0$.
From above, we have $\tilde a(s)= [a_s]$.
Note that $\tilde a$ is a loop, $a=a_1$ and $e_{z_0}=a_0$; therefore,
\[[a]=\tilde a(1)=\tilde a(0)=[e_{z_0}].\]
In other words $a\sim e_{z_0}$.
By \ref{sec:Homotopy of paths}, the corresponding homotopy can be lifted to $\tilde{\c{Z}}$ and it contracts $\tilde a$.
That is, any loop in $\tilde{\c{Z}}$ is contractable; thus, $\tilde{\c{Z}}$ is simply-connected.
\qeds

Recall that topological groups were defined in \ref{ex:group}.

\begin{thm}{Exercise}\label{ex:group-unicover}
Let $u\:\tilde{\c{G}}\to \c{G}$ be the universal covering of a reasonable topological group $\c{G}$
Show that $\tilde{\c{G}}$ admits a structure of topological group such that $u$ becomes a homomorphism of groups; that is $u(\tilde g\cdot \tilde h)=u(\tilde g)\cdot u(\tilde h)$ for any $\tilde g$ and $\tilde h$ in $\tilde{\c{G}}$.
\end{thm}

\section{Action of the fundamental group}

\begin{thm}{Theorem}\label{thm:action-on-ucovering}
Let $u\:(\tilde{\c{Z}},\tilde z_0)\z\to (\c{Z},z_0)$ be a pointed covering.
Suppose $\c{Z}$ is reasonable and $\tilde{\c{Z}}$ is simply-connected.
Then $u$ is a normal $\pi_1(\c{Z},z_0)$-covering.
\end{thm}

The following proof is similar to the construction of the monodromy homomorphism in \ref{sec:monodromy}, which in turn reuses the argument in \ref{thm:pi1S1Z}.

\parit{Proof.}
Choose a point $z_0\in\c{Z}$.
Let us denote by $\alpha\mapsto \tilde z_\alpha$ the bijection $\pi_1(\c{Z},z_0)\z\leftrightarrow u^{-1}\{z_0\}$ provded by \ref{thm:bijection-fiber}.

Let us recall, its construction.
Choose a point $\tilde z_1\in u^{-1}\{z_0\}$.
Given $\alpha=[a]\in \pi_1(\c{Z},z_0)$,
lift the loop $a\:[0,1]\to\c{Z}$ to a path $\tilde a\:[0,1]\to \tilde{\c{Z}}$ such that $\tilde a(0)=\tilde z_1$ and set
$\tilde z_\alpha=\tilde a(1)$.

By \ref{thm:lifting-criterion}, given $\alpha\in \pi_1(\c{Z},z_0)$ there is a unique continuous map $h_\alpha\:\tilde{\c{Z}}\to \tilde{\c{Z}}$ such that $h_\alpha\:\tilde z_1\mapsto \tilde z_\alpha$ and $u\circ h_\alpha=u$.

Given two elements $\alpha=[a]$ and $\beta=[b]$ in $\pi_1(\c{Z},z_0)$,
let $\tilde a$ and $\tilde b$ be the lifts of $a$ and $b$ that start at~$\tilde z_1$.
Note that $h_\alpha\circ \tilde b$ is a lift of $b$ that starts at $\tilde a(1)=\tilde z_\alpha$.
Therefore $\tilde a*(h_\alpha\circ \tilde b)$ is a lift of $a*b$ that start at $\tilde z_1$.
It follows that
\[h_{\alpha\cdot \beta}(\tilde z_1)=\tilde a*(h_\alpha\circ \tilde b)(1)=h_\alpha\circ\tilde b(1)=h_\alpha\circ h_\beta(\tilde z_1),\] and from uniqueness of lift, we get $h_{\alpha\cdot \beta}(\tilde z)=h_\alpha\circ h_\beta(\tilde z)$ for \textit{any} $\tilde z\in\tilde{\c{Z}}$.
It follows that $\alpha\mapsto h_\alpha$ is a homomorphism from $\pi_1(\c{Z},z_0)$ to the group of homeomorphisms of $\c{Z}$;
In other words $\alpha\mapsto h_\alpha$ defines a continuous action $\pi_1(\c{Z},z_0)\acts\tilde{\c{Z}}$.


Since $u$ is surjective and open, it is a quotient map.
The fibers of $u$ are exactly the orbits in $\tilde{\c{Z}}$.
Therefore $\c{Z}$ comes with the quotient topology of $\tilde{\c{Z}}$ by the constructed action, which proves the statement.
\qeds

\begin{thm}{Exercise}\label{ex:given-monodromy}
Let $\c{Z}$ be a reasonable connected space with base point $z_0\in\c{Z}$.
Given a group homomorphism $m\:\pi_1(\c{Z},z_0)\to G$, construct a normal $G$-covering of $\c{Y}\to\c{Z}$ with monodromy homomorphism~$m$.
\end{thm}


\section{From covering to subgroup}


\begin{wrapfigure}[5]{r}{48 mm}
\vskip-6mm
\centering
\begin{tikzcd}[column sep={1em}, row sep={2em}]
(\c{X},x_0)\arrow[dr, "v"']\arrow[rr, "h", dashed,leftrightarrow]& &(\c{Y},y_0)\arrow[dl, "w"]\\
&(\c{Z},z_0)&
\end{tikzcd}
\end{wrapfigure}

Let
$v\:(\c{X},x_0)\z\to (\c{Z},z_0)$ and $w\:(\c{Y},y_0)\z\to (\c{Z},z_0)$ be two pointed covering with connected reasonable base $\c{Z}$.
We say that they are \index{equivalent coverings}\emph{equivalent}, if there is a pointed homeomorphism $h\:(\c{X},x_0)\z\to(\c{Y},y_0)$ that makes the diagram commute.

\begin{thm}{Proposition}\label{prop:classification-coverings}
Let $\c{Z}$ be a connected reasonable space.
Then a pointed covering $v\:(\c{X},x_0)\z\to (\c{Z},z_0)$ is uniquely defined up to equivalence by the subgroup $H=\Im v_*<\pi_1(\c{Z},z_0)$.
\end{thm}

\parit{Proof.}
Let $G=\pi_1(\c{Z}, z_0)$.
By \ref{ex:induced-hom-covering:injective} $v_*\:\pi_1(\c{X},x_0)\z\to \pi_1(\c{Z},z_0)$ is a monomorphism; in particular, $H\cong \pi_1(\c{X},x_0)$.

\begin{wrapfigure}{r}{43 mm}
\vskip-2mm
\centering
\begin{tikzcd}
& (\c{X},x_0)\arrow[d, "v"]\\
(\tilde{\c{Z}},\tilde z_0) \arrow[r, "u"]\arrow[ur, "w"] & (\c{Z},z_0)
\end{tikzcd}
\end{wrapfigure}

Let $u\:(\tilde{\c{Z}},\tilde z_0)\to (\c{Z}, z_0)$ be the universal covering.
By \ref{thm:lifting-criterion} it can be included into the commutative diagram of coverings.
By \ref{thm:action-on-ucovering}, $u$ is a normal $G$-covering and $w$ is a normal $H$-covering.

Choose $\alpha=[a]\z\in \pi_1(\c{X},x_0)$.
Observe that the monodromy homeomorphism of $\alpha$ along $w$ is identical to the monodromy  homeomorphism of $v_*\alpha=[v\circ a]\in \pi_1(\c{Z},z_0)$ along $u$.

Indeed, let $\tilde a$ be a unique lift of $a$ along $w$ such that $\tilde a(0)=\tilde z_0$.
Note that $\tilde a$ is also a lift of $v\circ a$ along $u$.
Let $\tilde z_1=\tilde a(1)$, and let $h$ be the monodromy homeomorphism of $\alpha$ with respect to $w$.
By \ref{prop:monodromy-homo}, $h$ is uniquely defined by $w\circ h=w$ and $h(\tilde z_0) = \tilde z_1$ and it fits definition of monodromy homeomorphism of $v_*\alpha$ along $u$ as well.

We proved the

By \ref{thm:action-on-ucovering}, the space $\c{X}$ can be indentified with the quotient space $\tilde{\c{Z}}/H$ and the covering $v$ becomes a quotient map.
It follows that, the covering $w$ is equivalent to the covering $\tilde{\c{Z}}/H\to \tilde{\c{Z}}/G=\c{Z}$.

Recall that the induced homomorphism of composition is compostion of induced homomorphisms; see \ref{sec:Induced homomorphism}.
Therefore, we have the following commutative diagram of fundamental groups,
and the only-if part follows.

\begin{figure}[ht!]
\vskip-0mm
\centering
\begin{tikzcd}[column sep={1em}, row sep={2em}]
\pi_1(\c{X},x_0)\arrow[dr, "v_*"']\arrow[rr, "h_*", dashed,leftrightarrow]& &\pi_1(\c{Y},y_0)\arrow[dl, "w_*"]\\
&\pi_1(\c{Z},z_0)&
\end{tikzcd}
\end{figure}
\qeds

Denote by $m_\alpha$ and $n_\beta$ the monodromy
\ref{prop:monodromy-homo}


\section{An application}

%???PIC

A \emph{graph} is a topological space built from a discrete set of points called vertices and copies of the interval $[0,1]$ called edges, where each edge is attached by identifying its two endpoints with vertices.

Groups that isomorphic to the fundamental groups of graphs are called \index{free group}\emph{free}.

Usually free groups are defined in purely algebraic terms:
a group $G$ is free if there is a subset $S\subset G$ such that for any other group $H$, any map $S\to H$ can be extended to a unique homomorphism $G\to H$.
In this case, elements of $S$ are called \emph{generators} of $G$ and if $S$ and we say that a free group $G$ is generated by $S$.
By the extension property $G$ is defined by $S$ up to isomorphism.
Aside from the trivial group (which is a free group with 0 generators) the generating set of a free group is not uniquely defined, but the sets of generators have the same cardinality; so the term free group with $n$ generators makes sence.

There is also a combinatoric definition, where we start with the set $S$ and consider an alphabet that contains two letters $s$ and $s^{-1}$ for every $s\in S$; then
the set $G$ is identified with all finite words (including the empty word for neutral element) from this alphabet; the multiplication is defined as by combing two words in one with possible cancelations that only come from the definition of group.
Here is the beginining of the list of distinct elements of a free group with two generators $a$ and $b$; $1$ stands for the empty word:
\begin{align*}
&1,\\
&a, b, a^{-1}, b^{-1},\\
&a^2, a\cdot b, a\cdot b^{-1},
b\cdot a, b^2,b\cdot a^{-1},
a^{-1}\cdot b,a^{-2}, a^{-1}\cdot b^{-1},
b^{-1}\cdot a, b^{-1}\cdot a^{-1}, b^{-2},\\
&
a^3, a^2\cdot b, a^2\cdot b^{-1}, a\cdot b\cdot a,a\cdot b^{2},a\cdot b^{-1}\cdot a,  \dots
\end{align*}


One way or another, free groups are important in group theory and the following theorem, if proved in purely algebraic way, is not trivial, but as we are about to see, topological proof takes a few lines.
Proving equivalence of the definitions of free groups requires work, and the proof below should include it,
but I skip this part --- please forgive me this cheating.

\begin{thm}{Theorem}
Any subgroup of a free group is free.
\end{thm}

\parit{Proof.}
Choose a free group $F=\pi_1(\Gamma,v_0)$, where $\Gamma$ is a graph and $v_0$ is its vertex.
Note that graph is a reasonable space.
By \ref{prop:classification-coverings}, for any subgroup $H<F$ there is a covering $w\:\tilde\Gamma\to \Gamma$ and a vertex $\tilde v_0$ such that $w(\tilde v_0)=v_0$ and $\Im w_*=H$.
Recall that $w_*$ is a monomorphism, thus $\pi_1(\tilde \Gamma,\tilde v_0)$ is isomorphic to $H$.

Observe that the total space of a covering of graph is a graph.
Hence the statement follows.
\qeds

This statement gives some taste of possible applications of coverings in the group theory.
The following statements lead to more applications of that type.
\begin{itemize}
\item Any group is isomorphic to a group of some reasonable topological space.
\item Let $\Gamma$ is a connected graph and $T$ is a spanning tree in $\Gamma$.
Denote by $E$ the set of edges of $\Gamma$ that do not lie in $T$.
Then $\pi_1(\Gamma)$ is isomorphic to $F_E$.
\end{itemize}

\begin{thm}{Exercise}\label{ex:subgroup(free)}
Show that a nontrivial normal subgroup of infinite index in a finitely generated free group is not finitely generated.
\end{thm}

For more on the subject, read \cite{hatcher-2002} and check the references therein.

Let me add a text about subgroups and its translation into topological language.
Suppose $w\:(\c{Y},y_0)\to (\c{Z},z_0)$ be a pointed connected covering; assume that $\c{Z}$ is reasonable.
Suppose $G=\pi_1(\c{Z},z_0)$ and $H=\Im w_*\cong \pi_1(\c{Y},y_0)$.

\spell{\begin{multicols}{2}}{}

Recall that $H$ has index $n$ if $G$ has $n$ is cosets $g\cdot H$.

Furthermore, $H$ is a normal if the $g\cdot H=H\cdot g$ for any $g\in G$, in this case we have a natural homomorphism $G\to L=G/H$ and the quotient homomorphism $G\to L$ has kernel $H$.

\columnbreak

Recall that $w$ is $n$-fold if $w^{-1}(z_0)$ has exactly $n$ points.

Furthermore, the covering $w$ is normal if it is defined as a quotient map for a freely discontinuos action $L\acts\c{Y}$.
In this case, $H$ is normal, $L$ is isomorphic to $G/H$ and  have a monodromy homomorphism $G\to L$.

\spell{\end{multicols}}{}

Now assume we have two pointed connected covering $w_1\:(\c{Y}_1,y_1)\to (\c{Z},z_0)$ and $w_2\:(\c{Y}_2,y_2)\to (\c{Z},z_0)$; let $H_1=\Im w_{1*}\cong \pi_1(\c{Y}_1,y_1)$ and $H_2=\Im w_{2*}\cong \pi_1(\c{Y}_2,y_2)$.

\spell{\begin{multicols}{2}}{}
The intersection $K=H_1\cap H_2$ is a subgrop of $G$.
\columnbreak
The pointed fibered product of $w_1$ and $w_2$ has fundamental group isomorphic to  $K=H_1\cap H_2$.
\spell{\end{multicols}}{}


normalizator --- group of deck trasformations

intersection of subgropus --- fibered product of coverings restricted to a connected component.
